\section{Final Synthesis}
\label{sec:discussion:synthesis}

% The broader conclusion of this thesis is that robotic manipulation should move beyond a purely geometric view of action.
% Dynamic and non-prehensile interactions are not edge cases at the periphery of an otherwise geometric problem.
% They expose fundamental limitations in how robotic systems are designed, how manipulation skills are represented, and how learning-based methods are evaluated.

% The next generation of manipulation systems will likely require coupled reasoning over geometry, timing, contact, and dynamics rather than sequential pipelines that treat each in isolation.
% It will require better abstractions for skills and constraints than the atomic-skill libraries and ad hoc constraint encodings used today.
% It will require hybrid system architectures that combine learning with structured planning, so that strong priors and explicit feasibility refinement can coexist.
% It will require data collection methods that account for embodiment and dynamics, rather than only end-effector trajectories.
% And it will require honest evaluation of what statistical learning can and cannot generalize beyond its training distribution.

% Dynamics complicate manipulation, but they also expand what robots can do.
% The central opportunity is to design robotic systems that treat dynamics not as a disturbance to be avoided, but as a resource to be modeled, planned with, and exploited.

The broader conclusion of this thesis is that robotic manipulation cannot be fully understood through a purely geometric view of action.
Geometric reasoning remains essential: robots must still satisfy reachability, collision avoidance, grasp validity, and configuration-space constraints.
However, many of the behaviors that expand manipulation capability are governed by quantities that geometry alone does not capture.
Contact, inertia, timing, actuation limits, payload variation, fluid motion, and momentum transfer can determine not only whether a motion succeeds, but what kinds of motions are available in the first place.
Dynamic and non-prehensile interactions are therefore not edge cases at the boundary of manipulation.
They expose fundamental questions about how robotic systems represent feasibility, how skills are defined, how planners and learned models should interact, and how physical generalization should be evaluated.

The systems developed in this thesis address these questions from several directions.
Simulation-in-the-loop planning shows that dynamic interactions such as poking and whole-body contact can expose useful behaviors that are difficult to specify geometrically.
Reusable kinodynamic structure shows that dynamic feasibility can be amortized across related planning queries rather than recomputed from scratch.
Conditional generative models show that learned trajectory distributions can absorb information about payloads, failures, and constraint manifolds, enabling fast motion generation when the relevant physical parameters are made explicit.
Together, these results suggest that the next generation of manipulation systems should not treat geometry, timing, contact, dynamics, and learning as isolated components in a sequential pipeline.
They should instead expose these factors through interfaces that allow planners, optimizers, simulators, and learned models to share responsibility for physical admissibility.

This view also clarifies several requirements for future work.
Manipulation systems will need better abstractions for skills: abstractions that can represent not only grasping and transport, but also transitions among pushing, poking, sliding, caging, throwing, whole-body contact, and other dynamic interactions.
They will need better representations of constraints, so that dynamic feasibility can be searched, optimized, learned, reused, and checked at the appropriate level of the system.
They will need data collection methods that preserve the embodiment-dependent quantities that matter for dynamic behavior, rather than reducing demonstrations to end-effector trajectories alone.
They will also need evaluation protocols that distinguish interpolation within a training distribution from genuine generalization across new constraints, embodiments, and physical regimes.

Dynamics complicate manipulation because they make feasibility depend on motion history, contact conditions, timing, and the physical properties of both the robot and the world.
But those same dependencies expand the space of possible behaviors.
They allow robots to move objects without grasping them, exploit momentum rather than avoid it, adapt to altered embodiments, and operate in regimes where geometric pick-and-place is inefficient or impossible.
The central opportunity is therefore to design robotic systems that treat dynamics not as a disturbance to be suppressed after planning, but as structure to be represented, planned with, learned from, and exploited.
